\section{High Frequency Surveys}



Our setting is that of high-frequency repeated cross-sectional survey data. These surveys record time-specific output variables and time-independent covariates at the individual level. Such an array of variables would lend itself well to a fixed effects or mixed effects model were it not for the high frequency nature of the surveying. Such a frequency opens up the very likely possibility that noise correlated across time, breaking the key assumption of residual independence in standard fixed effects modeling. 

The mVAM survey responses comprise repeated cross-sectional (RCS) data, in which participants are not assumed to overlap from one day to the next. This distinguishes RCS surveys from panel data, in which a sample of respondents is fixed and repeatedly surveyed. Despite assumed independence in the selection of daily RCS samples, indicators collected through repeated surveys still retain a subtle, but nonetheless present, temporal correlation due to external factors~\citep{Lebo}.

Repeated cross-sectional data, as opposed to panel data, has less of an explicit correlation across time \citep{lebo2015repeated}. A number of existing of methods have arisen to account for this correlation, however none are well suited to surveys with such high frequency, individual-level covariates, and no time-dependent macrolevel variables. We combine elements from a multilevel model with a latent time series with a method for analyzing time series with added noise, generalizing both methods to enable their compatibility. The final product is an algorithm for fitting data from Equations \ref{eq:model_level1}-\ref{eq:model_level2}
